\chapter{Conclusiones}

Por todo lo investigado acerca del Teorema~\ref{thm:johnson-lindenstrauss}
se ha comprobado que es de gran interés en Ciencia de Datos y Machine Learning.
Los resultados están basados en el comportamiento de las matrices aleatorias como
método de proyección de un conjunto de puntos a una dimensión menor conservando
de forma aproximada las normas y las distancia entre los puntos.

La dimensión del espacio de proyección es del orden del
tamaño del conjunto de puntos e independiente de la dimensión original.


Estas buenas características han hecho del Teorema de Johnson-Lindenstrauss
una herramienta fundamental en técnicas de reducción de la dimensión,
especialmente en contextos donde se manejan grandes volúmenes de datos.


Su aplicación permite, en muchas ocasiones, en las mejoras y la
eficiencia computacional sin pérdida de información geométrica relevante,
lo que lo convierte en un pilar teórico clave en algoritmos de aprendizaje
automático, en la recuperación de información y en el análisis de datos
de alta dimensión.
